\scanpage{108}
\uline{Ex}\quad \(N=3\)

\textcircled{1}\quad \(\boxed{A^2(X)\ \text{fini}}\) en particulier \(J^2(X)=0\)

\textcircled{2}\quad
\diagram{\begin{array}{|l}\hline J^1(X)\longrightarrow J^1(Y)\quad\text{injectif exc. }p\\\hline\end{array}}
\deleted{\unsure{facteurs}}

\textcircled{3}\quad \deleted{\(\unsure{A^1(X)}\)} contenu dans 1°)

groupes essentiels \(J^1(X),N^1(X)\) définissant \(A^1(X)\)

hom essentiels
\diagram{\begin{array}{ll|l}
J^1(X)\longrightarrow J^1(Y)&&A^1(X)'\longrightarrow J^1(Y)\\
N^1(X)'\longrightarrow J^1(Y)&&
\end{array}}
\margin{\unsure{ssgrp} d'indice fini de \(N^1(X)\)}
\figure{\begin{tikzpicture}\draw(0,0)..controls(.5,.15)and(.75,-.15)..(1,.05);\end{tikzpicture}}

\uline{Pb} définir un complémentaire naturel de \(J^1(X)\) dans \(J^1(Y)\) (\unsure{correspdt.} aux ``cycles évanescents'')

Donner un sens à la formule
\[
J^1(X)\xrightarrow{\sim}J^1(Y)^\pi
\]
Interpréter le th. de \inserted{l'index de} Hodge pour les surfaces projectives en termes de la négativité d'une certaine forme définie sur \(N^1(X)\)\ldots ?

\uline{Ex}\quad \(N=4\)

\textcircled{1}\quad
\[
\left\{\begin{array}{l}N^2(X)\ \text{fini}\\A^3(X)\ \text{fini, en particulier }J^3(X)=0\end{array}\right.
\]
\textcircled{2}\quad
\diagram{\begin{array}{ll}
J^1(X)\longrightarrow J^1(Y)&\text{bij exc. }p\\
\deleted{J}N^1(X)\longrightarrow N^1(Y)&\text{inj. exc. }p\\
\Rightarrow\ A^1(X)\longrightarrow A^1(Y)&\text{\uline{bij exc.} }p
\end{array}}
\textcircled{\unsure{3}}
\margin{107}
